\section{Zeckendorf--窗口折叠：从微观态到稳定类型的可计算压缩}

\subsection{Zeckendorf 表示与窗口截断}

每个非负整数 $N$ 可唯一写为非相邻斐波那契数之和：
\[
N=\sum_{k\ge 1}c_k F_{k+1},\quad c_k\in\{0,1\},\quad c_kc_{k+1}=0.
\]
定义窗口折叠映射
\[
\mathrm{Fold}_m(N):=(c_1,\dots,c_m)\in X_m,
\]
必要时以 $0$ 补齐。

\subsection{可计算压缩与双重结构：稳定类型与退化度}

令 $\Omega_m=\{0,1,\dots,2^m-1\}$ 表示 $m$ 比特微观索引集合（以二进制解释），定义折叠映射
\[
\mathrm{Fold}_m:\Omega_m\to X_m.
\]
其将 $2^m$ 个微观态压缩为 $F_{m+2}$ 个稳定类型。对低端区间 $\{0,1,\dots,F_{m+2}-1\}$，$\mathrm{Fold}_m$ 为双射；对全区间 $\Omega_m$，$\mathrm{Fold}_m$ 为满射且存在非平凡退化度分布。
该压缩是算法性的：给定 $N$ 先求 Zeckendorf digits，再截断即可。

\subsection{\texorpdfstring{$m=6$}{m=6} 的最小刚性例}

当 $m=6$ 时，$\card{X_6}=F_8=21$，因此获得 $64\to 21$ 的最小压缩。并且存在循环/边界分裂：
\[
\card{X_6^{\mathrm{cyc}}}=18,\quad \card{X_6^{\mathrm{bdry}}}=3,
\]
其中边界词恰为
\[
X_6^{\mathrm{bdry}}=\{100001,\ 100101,\ 101001\}.
\]
此外，边界词在 $\mathrm{Fold}_6$ 下具有两重前像结构：每个边界词恰对应两个微观索引，其差为 $F_9=34$，形成“同一边界、不同提升量”的刚性几何。
上述结构为纯数学结论，且可由有限枚举完全复核，因此为模型的基础可审计输出。

